% !TEX root = ../main_final_en.tex
\section{Objectives and Justification}

The main objective of this study is to develop an artificial intelligence model capable of accurately predicting the capacity of Spanish forest plantations to absorb carbon dioxide (\(CO_2\)). This model is based on variables describing tree species, terrain characteristics and climatic and meteorological conditions. From this general objective, several specific goals are derived, which together justify the relevance and applicability of the project.

\subsection{Specific objectives}

The main objectives are listed below:

\begin{itemize}
    \item \textbf{Develop a robust predictive model:} Build a machine learning model that estimates the amount of \(CO_2\) that will be captured over time by a forest plantation, based on data such as species, soil type, diameter class, climate and other relevant variables.

    % \item \textbf{Optimise carbon capture:} Use the model to identify optimal combinations of species and terrain that maximise carbon fixation, contributing to the efficient planning of (re)forestation projects.

    \item \textbf{Ensure compliance with international regulations:} Guarantee that the model's predictions and outputs are compatible with the regulatory frameworks defined by the \textit{United Nations Framework Convention on Climate Change} (UNFCCC) and the \textit{Kyoto Protocol}, thus meeting the criteria required for the validation of carbon credits.

    \item \textbf{Analyse the determining factors of forest development:} Study the influence of climatic variables (such as temperature and precipitation) and edaphic variables (such as soil type or slope) on forest growth and its carbon capture capacity.

    \item \textbf{Support environmental and business decision-making:} Provide a practical and validated tool that enables technicians, managers and companies to select the most suitable species and plan afforestation actions with maximum efficiency in terms of carbon sequestration.
\end{itemize}

\subsection{Justification}

The need for predictive tools to estimate \(CO_2\) capture has intensified given the growth of the voluntary carbon credit market, and the obligations acquired: each country must report its greenhouse gas emissions and absorptions, and can use (re)forestation activities as offsetting mechanisms.

For these projects to be eligible, they must meet specific criteria, which make it essential to have models that not only estimate current carbon but are capable of forecasting its future evolution based on initial conditions and predictor variables.

This work seeks to fill that gap through the use of artificial intelligence applied to real, multi-source data. Integrating its management within the carbon credit system can represent an important opportunity for the local economy, which could support those regions where depopulation and the demographic challenge are a serious problem, and for the mitigation of climate change.


